
\medskip
\noindent\SGARefTarget{sga3:xi:proposition:3:12}\textbf{Proposition 3.12.} ---
\textit{Let \(S\) be a prescheme and \(G\) an affine \(S\)-group
prescheme.}

\begin{enumerate}
\item[\SGARefTarget{sga3:xi:proposition:3:12:item:a}\textit{(a)}]
\textit{Let
\[
  F:(\mathbf{Sch})^\circ/S\longrightarrow(\mathbf{Ens})
\]
be the functor such that, for every \(T\) over \(S\),}
\SGARefTarget{sga3:xi:display:3:12:finite-multiplicative-subgroups}
\[
  F(T)=
  \left\{
  \begin{array}{c}
  \text{multiplicative-type subgroups of \(G_T\)}\\
  \text{that are finite over \(T\)}
  \end{array}
  \right\}.
\]
\textit{Suppose that \(S\) is locally noetherian or that \(G\) is of
finite presentation over \(S\).  Then \(F\) is representable and
affine over \(S\).  If \(G\) is of finite presentation over \(S\),
then \(F\) is locally of finite presentation over \(S\).}

\item[\SGARefTarget{sga3:xi:proposition:3:12:item:b}\textit{(b)}]
\textit{Let \(H\) be a finite \(S\)-group prescheme of multiplicative
type.  Then
\(\operatorname{Hom}_{S\text{-}\mathrm{gr}}(H,G)\) is representable
and affine over \(S\).  If \(G\) is of finite type (respectively, of
finite presentation) over \(S\), then the same is true of
\(\operatorname{Hom}_{S\text{-}\mathrm{gr}}(H,G)\).}
\end{enumerate}

\medskip
\noindent\SGARefTarget{sga3:xi:remark:3:13}\textbf{Remark 3.13.} ---
Apart from the assertion that
\(\operatorname{Hom}_{S\text{-}\mathrm{gr}}\) is affine, and from the
case in which \(G\) is of finite presentation over \(S\), which is
the case that will suffice for us,
\SGARefLink{sga3:xi:proposition:3:12}{3.12} is an immediate
consequence of the theory of Hilbert schemes (A.~Grothendieck,
\textit{Techniques de construction et théorèmes d'existence en
Géométrie Algébrique: IV. Les Schémas de Hilbert}, Séminaire
Bourbaki, May 1961, no.~221).  It is enough, in fact, that \(G\) be
quasi-projective over \(S\).  In case \textit{\SGARefLink{sga3:xi:proposition:3:12:item:a}{(a)}} one can also
represent the larger functor
\[
  F_0(T)=
  \left\{
  \begin{array}{c}
  \text{group subpreschemes of \(G_T\) that are flat, proper,}\\
  \text{and of finite presentation over \(T\)}
  \end{array}
  \right\}.
\]
The canonical monomorphism \(F\to F_0\) is an ``open immersion,'' as
follows from the criterion in \SGARefLink{sga3:xi:proposition:3:1}{3.1} and \SGARefLink{sga3:x:corollary:4-7}{Exposé X, 4.7}~\textit{(b)}.
Thus the representability of \(F_0\) implies that \(F\) is
representable by an open subprescheme of \(F_0\).  In case \textit{\SGARefLink{sga3:xi:proposition:3:12:item:b}{(b)}}
it is enough to assume that \(H\) is projective and of finite
presentation over \(S\).  In both cases one obtains a functor locally
of finite presentation over \(S\).

In the present Exposé, \SGARefLink{sga3:xi:proposition:3:12}{3.12} is
only a technical lemma used to prove a key result in the \SGARefLink{sga3:xi:section:4}{next
section}.  We shall therefore sketch an easy direct proof that does
not use Hilbert schemes.

We first prove \textit{\SGARefLink{sga3:xi:proposition:3:12:item:b}{(b)}}.  It is enough to prove that
\(\operatorname{Hom}_S(H,G)\) is representable, independently of any
group structure on \(H\) and \(G\), and that it has the additional
properties asserted for
\(\operatorname{Hom}_{S\text{-}\mathrm{gr}}(H,G)\).  Indeed, applying
the same result to \(\operatorname{Hom}_S(H\times_S H,G)\), the
subfunctor
\(\operatorname{Hom}_{S\text{-}\mathrm{gr}}(H,G)\) of
\(\operatorname{Hom}_S(H,G)\) can be expressed by a finite inverse
limit---in fact, by fiber products---involving \(G\),
\(\operatorname{Hom}_S(H,G)\), and
\(\operatorname{Hom}_S(H\times_S H,G)\).  We leave this expression
to the reader.

Write
\[
  H=\operatorname{Spec}(B),
\]
where \(B\) is a sheaf of algebras on \(S\) that is locally free as a
sheaf of modules; this is the only hypothesis on \(H\) that the
argument uses.  Since the representability question is local on
\(S\), we may suppose that \(S\) is affine, with ring \(A\).  We may
also write \(G=\operatorname{Spec}(C)\), where \(C\) is an
\(A\)-algebra.

When \(G=S[t]=G_0\), with \(t\) an indeterminate, the Hom functor is
simply
\[
  T\longmapsto\Gamma(T,B_T).
\]
Since \(B\) is locally free, this functor is represented by the
vector bundle \(V(B^\vee)\), where \(B^\vee\) is the dual module
sheaf.  When \(G=S[(t_i)_{i\in I}]=G_0^I\), for a not necessarily
finite family of indeterminates, it follows that\footnote{The source
display prints
\(\operatorname{Hom}_S(H,G)^I=V(B^\vee)^I=V(B^{\vee(I)})\).
Since the preceding sentence sets \(G=G_0^I\), the displayed equality
requires the explicit factor
\(\operatorname{Hom}_S(H,G)=\operatorname{Hom}_S(H,G_0)^I\).}
\SGARefTarget{sga3:xi:display:3:13:hom-affine-space}
\[
  \operatorname{Hom}_S(H,G)
  =\operatorname{Hom}_S(H,G_0)^I
  =V(B^\vee)^I
  =V\bigl(B^{\vee(I)}\bigr).
\]

In general \(G\) is isomorphic to a closed subscheme of a scheme of
the form \(S[(t_i)_{i\in I}]\); equivalently, \(C\) is a quotient of
an \(A\)-algebra \(A[(t_i)_{i\in I}]\).  Let \((F_j)\) be a family of
generators of the ideal by which one takes the quotient.  After
covering \(S\) by smaller affine open subsets, we may assume that
\(B\) is free of rank \(n\).  Choose a basis
\((e_k)_{1\leq k\leq n}\) of \(B\).  Writing
\[
  x_i=\sum_{k=1}^{n}x_{ik}e_k
\]
with indeterminate coefficients \(x_{ik}\) in an unspecified
\(A\)-algebra \(A'\), expand \(F_j((x_i))\) in this basis.  For every
\(F_j\) this gives \(n\) polynomials
\[
  F_{j,k}\qquad(1\leq k\leq n)
\]
in the variables
\((x_{i,k})_{i\in I,\,1\leq k\leq n}\), with coefficients in \(A\).
It follows at once that \(\operatorname{Hom}_S(H,G)\) is represented
by the spectrum of the quotient of the polynomial ring
\(A[(x_{i,k})]\) by the ideal generated by the \(F_{j,k}\).  This
proves \textit{\SGARefLink{sga3:xi:proposition:3:12:item:b}{(b)}}.

We now prove \textit{\SGARefLink{sga3:xi:proposition:3:12:item:a}{(a)}}.  For every ordinary finite commutative
group \(M\), let \(F_M\) be the subfunctor of \(F\) obtained by
restricting to those subgroups of \(G_T\) that are of multiplicative
type and of type \(M\).  A gluing argument like the one used in \SGARefLink{sga3:xi:proposition:3:5}{3.5}
(which we should have stated after decomposing the problem a little
further) shows that it is enough to prove that each \(F_M\) is
representable.  Then \(F\) is represented by the disjoint sum of the
\(F_M\), with \(M\) ranging over the isomorphism classes of finite
commutative groups.  Indeed, \(F\) is the corresponding sum in the
category of sheaves.

Fix \(M\) and write \(F\) in place of \(F_M\).  Put
\(H=D_S(M)\), and let
\[
  F'=\operatorname{Imm}_{S\text{-}\mathrm{gr}}(H,G)
  \subset\operatorname{Hom}_{S\text{-}\mathrm{gr}}(H,G)
\]
be the subfunctor whose value at \(T\) is the set of homomorphisms of
\(T\)-groups \(H_T\to G_T\) that are closed immersions.  By
\textit{\SGARefLink{sga3:xi:proposition:3:12:item:b}{(b)}},
\(\operatorname{Hom}_{S\text{-}\mathrm{gr}}(H,G)\) is represented by
an affine \(S\)-prescheme.  Using \SGARefLink{sga3:ix:theorem:6:8}{Exposé IX, 6.8}, one sees immediately
that \(F'\) is represented by an open-and-closed subprescheme of it;
in particular, \(F'\) is affine over \(S\).

Consider the canonical morphism \(F'\to F\), which assigns to every
monomorphism \(H_T\to G_T\) its image group.  Directly from the
definitions, \(F'\) is thereby a principal homogeneous bundle over
\(F\), in the category of fpqc sheaves, under
\[
  \Gamma_F=\Gamma_S\times_S F,
  \qquad
  \Gamma=\operatorname{Aut}_{\mathrm{gr}}(M).
\]
It follows by descent that \(F'\to F\) is representable: for every
morphism \(T\to F\) with \(T\) representable, the prescheme
\(T\times_F F'\) is representable over \(T\), in fact by a principal
Galois bundle under \(\Gamma\).

By \SGARefLink{sga3:iv:theorem:4:6:6}{Exposé IV, 4.6.6}, \(F\) is therefore representable if and only if
the quotient \(F'/F''\), where \(F''=F'\times_F F'\), exists in
\((\mathbf{Sch})\) and is universally effective for faithfully flat
quasi-compact morphisms.  Equivalently, the quotient \(F'/\Gamma\)
must exist and be universally effective for such morphisms.  Since
\(F'\) is affine over \(S\), \SGARefLink{thm:sga3-V-4.1}{Exposé V, 4.1} shows that this condition
is satisfied.  This proves the representability of \(F\) in
\textit{\SGARefLink{sga3:xi:proposition:3:12:item:a}{(a)}}.

Finally, suppose that \(G\) is of finite presentation over \(S\).
The additional assertion follows immediately from the preceding
proof: by \textit{\SGARefLink{sga3:xi:proposition:3:12:item:b}{(b)}}, \(F'\) is then locally of finite presentation
over \(S\), so that Exposé V may be applied.

